\section{Conclusion and Open Problems}
\label{sec:conclusion}

We prove several implications and non-implications between fairness notions,
and for additive valuations, we give an almost complete picture of implications.
We believe this would help inform further research in fair division.
This would be especially useful if one wants to extend a fair division result
to a stronger notion or a more general setting,
or study a weaker notion or a simpler setting for a hard problem.

\Cref{fig:additive-nny} shows us that the only notions whose feasibility is unknown are EFX and PMMS,
and resolving their feasibility is one of fair division's most important problems.
For mixed manna, even the existence of MXS allocations is open.

For equally-entitled agents having additive valuations over goods or over chores,
EF1+PO allocations are known to exist \cite{caragiannis2019unreasonable,barman2018finding,mahara2025existence},
but their efficient computation remains open.
Relaxing the problem to EEF1+PO can be a helpful first step.

Here are four interesting open problems regarding implications that we could not resolve:
\begin{tightenum}
\item For additive goods (equal entitlements), does MXS imply EEF1?
    Note that the implication holds for the special case of two agents
    (\cref{thm:impl:mxs-to-ef1-n2} in \cref{sec:impls-extra}).
\item For additive goods (unequal entitlements), does APS imply PROP1?
    This is open even for two agents.
\item For submodular goods (equal entitlements), does MXS imply PROP1?
    This is true for binary marginals
    (\cref{thm:impl:m1s-to-propx-binary-subadd} in \cref{sec:impls-extra}).
\item For submodular goods with binary marginals (equal entitlements),
    does pairwise-MMS imply MMS?
\end{tightenum}
For less-studied settings, like non-additive valuations over chores,
many implications are still open.

Another interesting direction is to study implications of the form $F_1$+PO $\fimplies$ $F_2$+PO.
For additive valuations and equal entitlements, most questions of this form are already resolved.
This is because if $F_1 \fimplies F_2$, then $F_1$+PO $\fimplies$ $F_2$+PO.
On the other hand, most of our counterexamples use identical valuations,
where every allocation is trivially PO.

Some fairness notions and settings don't fit our model (\cref{sec:prelims}),
so it is unclear how to systematically represent results about them
and extend the inference engine (\cref{sec:cpig}) for them.
Examples of such settings include constrained fair division
\cite{gourves2014near,bouveret2017fair,biswas2018fair,equbal2024fair},
notions based on social graphs \cite{aziz2018knowledge},
multiplicative approximations of fairness notions \cite{amanatidis2018comparing},
and parametrized notions like $\mathrm{WEF}(x, y)$ \cite{chakraborty2024weighted}.
Extending our techniques to these settings and notions would be an interesting line of research.

We omitted some fairness notions from our work because
they are fundamentally different from the notions we consider.
Equitability (EQ) \cite{brams1996fair}, and its relaxations like EQ1 and EQX
\cite{amanatidis2023fair,gourves2014near}, compare different utility functions with each other.
Notions like CEEI \cite{varian1974equity} and maximum Nash welfare \cite{caragiannis2019unreasonable}
include an aspect of efficiency in addition to fairness.
We also didn't study fair division of divisible items \cite{steinhaus1940sur,stromquist1980how,varian1974equity},
or a mix of divisible and indivisible items \cite{liu2024mixed}.
%Another reason to not study these notions and settings was to limit the scope of our work,
%since our results are already quite extensive,
%(see \cref{table:impls1,table:cex-add,table:cex-nonadd,table:impls-tribool,table:impls-unit-demand}).
Nevertheless, we believe that these notions and settings offer an interesting line of research,
and our inference engine (\cref{sec:cpig}) can be readily adapted for them.
